\documentclass[11pt,a4paper]{article}
\usepackage[utf8]{inputenc}
\usepackage[T1]{fontenc}
\usepackage[margin=2.4cm]{geometry}
\usepackage{parskip}
\usepackage{titlesec}
\usepackage{xcolor}
\usepackage{tcolorbox}
\usepackage{fancyhdr}

\tcbuselibrary{skins,breakable}

\definecolor{nv}{RGB}{15,23,42}
\definecolor{bl}{RGB}{0,118,206}
\definecolor{mg}{RGB}{226,0,116}
\definecolor{mut}{RGB}{100,116,139}

\titleformat{\section}{\large\bfseries\color{nv}}{}{0em}{}
\titleformat{\subsection}{\normalsize\bfseries\color{bl}}{}{0em}{}

\pagestyle{fancy}
\fancyhf{}
\fancyhead[L]{\small\color{nv}Guion de presentación}
\fancyhead[R]{\small\color{nv}\thepage}
\renewcommand{\headrulewidth}{0.4pt}

\newtcolorbox{slide}[1]{colback=white,colframe=mg,boxrule=0.5pt,arc=2pt,
  left=10pt,right=10pt,top=8pt,bottom=8pt,breakable,
  title={\textbf{#1}},coltitle=white,colbacktitle=nv,fonttitle=\small\bfseries}

\title{\vspace{-1.2cm}\textbf{Guion — Generación autónoma de microservicios}\\[0.2cm]
\large Notas del ponente para el deck TCS (11 diapositivas)}
\author{}
\date{}

\begin{document}
\maketitle
\thispagestyle{fancy}

\begin{tcolorbox}[colback=bl!6,colframe=bl,arc=2pt,left=8pt,right=8pt,top=6pt,bottom=6pt]
\textbf{Duración total:} $\approx$ 4 minutos. Ritmo pausado; cada diapositiva se lee
en 20--25 segundos. Frases cortas, una idea por frase.
\end{tcolorbox}

\section*{Apertura}

\begin{slide}{Diapositiva 1 · Portada}
Buenos días. Les voy a presentar en pocos minutos cómo generamos microservicios de
forma autónoma: dónde está la industria, cómo funciona nuestra solución, cuánto
cuesta y hacia dónde vamos.
\end{slide}

\section*{El contexto}

\begin{slide}{Diapositiva 2 · Dónde está la industria}
Hay tres caminos en el mercado. Primero, los agentes generalistas que arreglan código
ya existente: dominan los rankings, pero están saturados y salen caros. Segundo, los
generadores guiados por especificación --- que construyen desde un diseño formal ---
que es nuestro camino, y donde todavía queda mucho espacio: solo se resuelve entre el
12 y el 55 por ciento. Y tercero, los agentes que se auto-mejoran: suenan muy bien,
pero en código real la mejora es casi nula. La lección es simple: lo que de verdad
marca la diferencia no es un modelo más listo, sino una \emph{verificación} en la que
puedas confiar.
\end{slide}

\section*{La solución}

\begin{slide}{Diapositiva 3 · Cómo funciona}
¿Cómo funciona? Entra una especificación: entidades, reglas, historias de usuario. Un
agente autónomo escribe el servicio completo en cinco etapas. Después, un contenedor
sellado y sin red lo compila y lo prueba. Si pasa, se entrega; si falla, el agente
corrige, hasta tres veces. Al tercer fallo se detiene y pide ayuda a una persona.
Nunca un bucle infinito, nunca un éxito no comprobado.
\end{slide}

\begin{slide}{Diapositiva 4 · El proceso, de punta a punta}
El proceso completo tiene nueve etapas: de los requisitos a la entrega, pasando por
arquitectura, modelos, generación, calidad y despliegue. La generación y la
verificación están en el centro.
\end{slide}

\begin{slide}{Diapositiva 5 · Lo que eso garantiza}
¿Qué garantiza esto? Un servicio completo, no fragmentos. Verificación en un entorno
aislado. Autocorrección con tope. Y reglas del proyecto que se auditan de forma
automática --- no la opinión del modelo sobre su propio trabajo.
\end{slide}

\section*{El coste}

\begin{slide}{Diapositiva 6 · El coste}
¿Cuánto cuesta? Un servicio de dificultad media: siete centavos de dólar. Un servicio
\emph{garantizado} --- generamos tres y entregamos el que compila --- veinte centavos.
El acumulado de todo el trabajo hasta hoy: ocho dólares con quince centavos.
\end{slide}

\begin{slide}{Diapositiva 7 · Coste por etapa}
Y el coste es predecible: la etapa más cara es la de tests; la más barata, la de
dominio. Escala con el tamaño del servicio.
\end{slide}

\begin{slide}{Diapositiva 8 · Tokens}
En tokens: diecinueve millones acumulados, unos setenta y ocho mil por servicio. Todo
medido y visible en MLflow.
\end{slide}

\section*{Balance y futuro}

\begin{slide}{Diapositiva 9 · El trabajo hecho}
¿Qué logramos hasta ahora? La plataforma ya verifica su propia salida --- antes
reportaba cero de quince verificadas, por un defecto de un carácter. Las dos vías de
producto compilan y prueban. La métrica de calidad dejó de ser manipulable. Y el
registro de coste es honesto. Queda un punto abierto: el verificador a veces cambia
su veredicto, y lo cerramos en la siguiente iteración.
\end{slide}

\begin{slide}{Diapositiva 10 · El futuro: auto-mejora}
¿Hacia dónde vamos? A que el sistema aprenda de sus errores y mejore su propio
manual. Pero la evidencia nos obliga a un orden: primero un verificador sólido;
después, una medida que detecte fallos reales; y solo entonces, encender la
auto-mejora. El mecanismo ya está construido.
\end{slide}

\begin{slide}{Diapositiva 11 · El camino}
En resumen: pasamos de ``no podía verificar su propia salida'' a ``genera, compila y
verifica por siete centavos, y se auto-mejorará cuando la base sea fiable''. El camino
son cuatro pasos cortos. Gracias --- y quedo para preguntas.
\end{slide}

\end{document}
